\section{Conclusions}
\label{sec:conclusions}

\subsection{What this paper establishes, in plain language}

Hexo gives a player two stone placements on an ordinary turn.  A line of six
wins as soon as its last stone is placed.  The board is unbounded, but a legal
placement must lie within distance eight of an existing stone.  These rules
create a search problem with a local tactical goal and a broad legal reply
set.  A midgame defender node may have hundreds of legal cells.  Checking all
of them at every level is often the main cost of a threat-space search.

The basic result is a safe way to omit some of those cells.  The omission is
not justified by distance alone.  It is justified by a finite certificate.
The certificate records the attacker's planned placements, the empties needed
by terminal threat witnesses, the number of defender opportunities before
each use, and the windows in which a hidden defender stone could complete a
line.  A verifier forms a zone from those records.  Direct obligations are
searched.  A future obligation that is not yet legal receives a seed band
wide enough to cover every possible radius-eight relay.  Defender-alive
windows receive separate guards according to whether they already contain a
defender stone or are still empty.

The dismissal theorem says that a certificate passing these checks survives
every legal defender reply, including replies outside the searched set.  Its
proof follows a real game and a ghost game together.  The real game contains
the omitted reply.  The ghost game contains a searched filler and follows the
certificate.  The zone prevents the first omitted stone from occupying a
future required cell or completing a defender window.  The attacker therefore
reaches the same certified resolution, or a permitted adaptive LOSS
resolution, by the declared placement horizon.  Search trimming can still
miss an attacker win.  It cannot turn a losing attack into an accepted win
when the verifier enforces the certificate conditions.

The zone can be computed with cell-specific deadlines and window-specific
exposures.  This is no larger than assigning every object the full certificate
depth, and can be smaller.  A tight current threat family admits a different reduction.  At a
protected forcing gate, an exact copied defender hit does not create a
real-only legality frontier.  The debited clocks remove that opportunity from
the relevant seed band and charge it to a selected window only when the hit
lies in that window.  A defender that refuses the gate is handled by an
explicit escape contract.  This result is conditional on exact copying and
the gate checkpoint.  It does not reduce the global scalar budget and does
not make a substituted hit free.

Terminal LOSS certificates also have a bounded description.  With one
defender placement left, three threat windows suffice.  With two placements
left, five suffice.  The triangle and five-cycle families show that these
caps are exact for the named-window witness format.  The result reduces proof
data without changing the universal survivor condition.

Two global defenses have exact boundaries.  A position-independent pairing
cannot cover every length-six window on the three-axis hex lattice.  At
length seven, a periodic perfect matching exists, covers every window exactly
once, and has minimum period index twelve among periodic constructions.  The
potential method gives sound fixed-family, finite-horizon, and finite-region
blocking certificates at base $\sqrt3$.  It does not give a global dynamic
greedy defense: an exact 124-branch enumeration gives a fixed attacker line
that defeats every greedy tie-breaking from a position with $\Phi<1$.  A
non-greedy global defense from every such position is not ruled out.

The resulting rule is concrete.  A solver may skip a reply only with a
checked reason tied to the selected proof.  Exact-copy gates, three/five LOSS
witnesses, fixed-family potential bounds, and the narrow domination patterns
provide such reasons under their premises.  Distance-only pruning,
unrestricted dynamic greedy, and a global pairing for length six do not.

\subsection{Results and limits}

\Cref{tab:conclusion-results} separates the result consumed by a verifier
from its main limit.  The detailed quantitative boundary is the limit map in
\cref{sec:tightness}.

\begin{table}[tbp]
  \centering
  \caption{Results and their operative limits.}
  \label{tab:conclusion-results}
  \small
  \renewcommand{\arraystretch}{1.1}
  \begin{tabular}{@{}p{.25\textwidth}p{.32\textwidth}p{.33\textwidth}@{}}
    \toprule
    Layer & Established result & Operative limit \\
    \midrule
    Dismissal & A valid zone certificate preserves the attacker win against
      every legal defender reply by its horizon. & The result is
      certificate-relative and may lose search completeness. \\
    Ranked zones & Exact role ranks and per-window exposures determine finite
      direct, seed, touched, and virgin terms. & Uniform-band and full-union
      virgin sharpness remain open. \\
    Forced-hit gates & Protected exact-copy gates permit $f$ and $Q^{\D}$
      hazard debits. & No scalar $B$ debit, free D17 substitution, automatic
      net shrinkage, or slack-pressure debit is proved. \\
    LOSS witnesses & Three windows suffice at $b=1$ and five at $b=2$. & The
      lower examples pin the named-window witness representation. \\
    Pairing & No global covering matching exists for $k=6$; a minimal-index
      periodic construction exists for $k=7$. & Position-dependent partial
      pairings are not excluded. \\
    Potential & Fixed-family and finite-region inequalities certify blocking;
      finite horizons admit explicit bounds. & Dynamic touched-window greedy
      fails globally; the raw non-greedy claim is open. \\
    Domination & P1, P2, and P3 permit checked local quotienting. & Their
      frontier, dead-spoke, initial-legality, and singleton-nonwin premises
      are mandatory. \\
    Machine checks & Pairing and the greedy counterexample have exhaustive
      finite enumerations; other runs test named calculations or samples. &
      Search experiments and randomized spot checks are not theorem proofs. \\
    \bottomrule
  \end{tabular}
\end{table}

Rule constants such as window length six, eighteen incident windows, and the
radius-eight legality step are not tuning parameters.  Within the certificate
format, $8(r-1)$, the LOSS deadline $p_{\rm leaf}+b+2$, the three/five witness
caps, and the leaf-entry and OR-COMPLETION deadlines have attained or absolute
boundaries.  Scalar path maxima, transition charges, and exact-kernel tests
are instead framework-relative; indexed fields or another grammar may encode
the same hazards differently.

The uniform $8(B-1)$ wrapper and full-union virgin radius remain open.  The
gate calculus is partial: it debits exact copied hits from named hazards, but
an escape, LOSS remainder, or substitution retains the full clock.  Added
checkpoint roles may offset a smaller radius.  Neither a fixed-window equality
nor a hazard debit promises a smaller total set at every node.

\subsection{Implications for solvers}

A df-pn implementation can separate search from verification.  Search may
use any move ordering or candidate generator.  It may start with the measured
radius-three and attacker-touched predecessor set, use threat incidence for
ordering, or learn a proposal distribution.  When that search finds a
candidate win, a compiler attaches the finite grammar, exact child map,
obligations, clocks, window data, and terminal witnesses.  An independent
verifier then checks the theorem selected at each defender node.  Failure of
one check yields an unverified or unknown result, not an attacker win.

At an ordinary node, the verifier recomputes legality and masks, then checks
protected cells, ranks, window exposures, searched-set containment, attacker
legality, terminal data, and the horizon.  A DAG also needs one consistent
position, phase, label, and clock assignment at each shared node.
Branch-indexed substitution uses its transition-inclusive envelope.

Deep certificates gain most from exact labels.  A uniform band uses the
largest remaining defender budget around every future carrier.  A ranked band
uses only the defender opportunities before that carrier's own deadline.
Windows that the attacker enters early stop their exposure clocks.  A DAG can
share repeated suffixes without increasing the selected path budget.  These
reductions target the part of the proof where a coarse horizon creates the
largest spatial bands.

Forcing lines admit an additional node type.  The verifier first checks the
named current threat masks, $\tau(F_Q)=b$, and the absence of an immediate
defender win.  It protects every threat empty through that checkpoint.  It
then searches exactly the extendable-hit kernel.  A kernel reply is copied at
the same cell and cannot use D17 substitution.  A reply outside the kernel
abandons the old subtree and invokes the finite escape contract with deadline
$p(Q)+b+2$.  Ordinary descendants may use the debited clocks.  A loose threat
family with $\tau(F_Q)<b$ does not qualify for this reduction.

At a LOSS leaf, the compiler may retain at most three or five windows.  The
verifier still checks their masks, completion legality, universal survivor
condition, and defender non-win diagnostic.  The cap does not validate a
heuristic threat list.

No fixed partner table covers length six.  Potential checks apply to their
fixed-family or finite-region premises; dynamic danger is otherwise an
ordering heuristic.  P1, P2, and P3 also retain their frontier, dead-spoke,
initial-legality, and singleton-nonwin checks.  Ordered intermediate nodes
remain when a certificate depends on them.

The measured 147-cell predecessor set against roughly 302 legal cells is a
generator target, not a theorem zone.  Learned proposals and node caps have
the same status.  They may shorten search, while only an accepted verifier
rule reduces the defender set in the proof.

\subsection{Future work}

\paragraph{Broader gate debits.}
The forced-hit result leaves four connected questions.  A target-independent
scalar budget smaller than $B$ would shrink every coarse wrapper, but it must
cover terminal ordering, LOSS remainders, and every window at once.  A
substituted hit can change the legality frontier even when it answers the same
threat, so a sound D17 debit needs a stronger domination or frontier-inertness
certificate.  Checkpoint roles may offset the saved hazard radius, and a
tight-gate theorem alone does not prove net searched-set shrinkage.  A family
with transversal number below the placement budget also leaves a quiet move;
a finite residual-state account for that slack could extend the gate method.
Resolving these questions would determine when forcing pressure produces a
predictable end-to-end branch reduction.

\paragraph{Uniform and full-union geometry.}
The uniform $8(B-1)$ obligation wrapper may contain slack not present in the
exact ranks.  A smaller bound would simplify certificates that do not carry
cell-specific deadlines.  It requires either a complete equality family or a
role-and-turn theorem showing that a ghost-illegal carrier always has less
than the full budget.  The full-union virgin radius has a different obstacle:
overlap among the eighteen incident windows can select a seed that the target
window no longer selects.  A solution would either construct certificates
excluding every competing term or prove that overlap always yields a smaller
safe union.

\paragraph{Potential boundary.}
Five questions remain distinct.  GAP-RAW asks whether every nonterminal
defender-FirstStone position with $\Phi<1$ has some non-greedy forever
defense.  GAP-GLOBAL-RENEWAL asks for an invariant that absorbs the stream of
new touched windows.  GAP-AMORTIZED-ABANDONMENT asks whether killed or
abandoned targets can fund that renewal without a divergent virgin sum.
GAP-DYNAMIC-PAIRING asks whether pairs chosen with the position can complement
the potential account even though a static global pairing is impossible.
GAP-FINITE-CUTOFF asks for an a priori horizon bound for a finite witness to
failure or survival.  A solution to the last question would turn the
K\"onig-style finite reduction into a bounded decision procedure.

\paragraph{Domination and local enumeration.}
The P1 and P2 premises are deliberately narrow.  A richer causal boundary
state could track unmatched frontier chains and live mask differences for a
fixed horizon.  Such a state may certify more dead-cell replacements or more
interchangeable hits.  It must preserve early terminality, actor-directed
legal inclusion, and every crossing window.  Running the full MV-P1, MV-P2,
and MV-P3 quotient specifications would add independent exhaustive checks of
the present local patterns.  It would not by itself prove a broader pattern.

\paragraph{Formalization.}
The certificate objects are finite.  Positions have finite occupancy,
windows queried by a certificate form finite families, transversal tests have
budgets at most two, and the tree or DAG is finite.  These features make a
Lean or Coq reference checker feasible.  A practical order is to formalize
axial distance and windows, the phase machine and terminal order, the D9
grammar, and the scalar recurrences.  The real--ghost coupling, ranked
first-bad-event proof, transition envelopes, and gate escape proof can then be
layered on those definitions.

The work is substantial.  The unbounded board requires careful finite-support
lemmas.  LOSS leaves quantify over a finite but adaptive defender remainder.
DAG sharing needs path clocks and label consistency.  Gate checkpoints add a
two-phase protected set and a separate escape horizon.  Potential boundary
results use exact algebraic comparisons and a compactness argument.  A
formalization would give a small trusted verifier and expose omitted scope
conditions, but it would not certify the search heuristics or settle the OPEN
rows without new mathematics.
